%!TeX program = xelatex
\documentclass{SYSUReport}
\usepackage{tabularx}
\usepackage{float}

% 根据个人情况修改
\headl{}
\headc{}
\headr{并行程序设计与算法实验}
\lessonTitle{并行程序设计与算法实验}
\reportTitle{Lab9-CUDA矩阵转置}
\stuname{xxx}
\stuid{xxx}
\inst{计算机学院}
\major{xxx}
\date{\today}

\begin{document}

% =============================================
% Part 1: 封面
% =============================================
\cover
\thispagestyle{empty}
\clearpage

\section{实验目的}
\begin{itemize}
   \item 熟悉 CUDA 线程层次结构(grid、block、thread)和观察 warp 调度行为。
   \item 掌握 CUDA 内存优化技术(共享内存、合并访问)。
   \item 理解线程块配置对性能的影响。
\end{itemize}

\section{实验内容}
\subsection{CUDA并行输出}
\begin{enumerate}
    \item 创建$n$个线程块，每个线程块的维度为$m\times k$。
    \item 每个线程均输出线程块编号、二维块内线程编号。例如：
    \begin{itemize}
        \item “Hello World from Thread (1, 2) in Block 10!”
        \item 主线程输出“Hello World from the host!”。
        \item 在 main 函数结束前，调用 \texttt{cudaDeviceSynchronize()}。
    \end{itemize}
    \item 完成上述内容，观察输出，并回答线程输出顺序是否有规律。
\end{enumerate}

\subsection{使用 CUDA 实现矩阵转置及优化}
\begin{enumerate}
    \item 使用 CUDA 完成并行矩阵转置。
    \item 随机生成 $N \times N$ 的矩阵 A。
    \item 对其进行转置得到 $A^T$。
    \item 分析不同线程块大小、矩阵规模、访存方式(全局内存访问，共享内存访问)、任务/数据划分和映射方式，对程序性能的影响。
    \item 实现并对比以下矩阵转置方法：
    \begin{itemize}
        \item 仅使用全局内存的 CUDA 矩阵转置。
        \item 使用共享内存的 CUDA 矩阵转置。
        \item 使用 padding 优化共享内存访问，减少存储体冲突。
    \end{itemize}
\end{enumerate}

\section{代码说明}
\subsection{CUDA Hello World}
\texttt{src/hello\_world.cu} 接收三个整数 $n,m,k$，分别表示线程块数和每个线程块的二维维度。程序在 host 端输出提示信息，然后启动 $n$ 个线程块，每个线程块包含 $m\times k$ 个线程，并调用 \texttt{cudaDeviceSynchronize()} 等待 device 端输出完成。

\subsection{CUDA 矩阵转置}
\texttt{src/matrix\_transpose.cu} 使用 CUDA event 计时，并对三种 kernel 分别计量平均运行时间：
\begin{itemize}
    \item \textbf{全局内存版本}: 每个线程直接读取 $A[row][col]$ 并写入 $A^T[col][row]$。
    \item \textbf{共享内存版本}: 每个线程块先把矩阵 tile 读入共享内存，再交换块坐标写回全局内存。
    \item \textbf{优化共享内存版本}: 共享内存 tile 的列数加 1，缓解转置读写时的 shared memory bank conflict。
\end{itemize}

\section{实验结果与分析}
\subsection{CUDA Hello World 并行输出}
\subsubsection{实验现象}
描述实验观察到的现象，例如线程输出的顺序等。可以粘贴部分关键的运行截图或输出文本。
\begin{itemize}
    \item 回答：
\end{itemize}
\subsubsection{结果分析}
线程输出顺序是否有规律？为什么？结合CUDA线程调度机制进行解释。
\begin{itemize}
    \item 回答：
\end{itemize}

\subsection{CUDA 矩阵转置及优化}
\subsubsection{不同实现方法的性能对比}
展示不同矩阵转置实现（仅全局内存、使用共享内存、优化共享内存访问）在不同矩阵规模 (N) 和不同线程块大小下的运行时间。

\begin{table}[h!]
\centering
\caption{矩阵转置性能对比 (时间单位: ms)}
\begin{tabular}{|c|c|c|c|c|}
\hline
矩阵规模 (N) & 线程块大小 & 全局内存版本 & 共享内存版本 & 优化共享内存版本 \\
\hline
\multirow{3}{*}{512} & 8$\times$8 & & & \\
\cline{2-5}
& 16$\times$16 & & & \\
\cline{2-5}
& 32$\times$32 & & & \\
\hline
\multirow{3}{*}{1024} & 8$\times$8 & & & \\
\cline{2-5}
& 16$\times$16 & & & \\
\cline{2-5}
& 32$\times$32 & & & \\
\hline
\multirow{3}{*}{2048} & 8$\times$8 & & & \\
\cline{2-5}
& 16$\times$16 & & & \\
\cline{2-5}
& 32$\times$32 & & & \\
\hline
\end{tabular}
\end{table}

\subsubsection{结果分析}
1. 根据实验结果，总结线程块大小、矩阵规模对程序性能的影响。哪种配置下性能最优？为什么？

回答：

2. 讨论任务/数据划分和映射方式对性能的影响。

回答：

注：实验报告格式参考本模板，可在此基础上进行修改；实验代码以zip格式另提交；最终提交内容包括实验报告(pdf格式)和实验代码(zip压缩包格式)
\end{document}
